% !TeX program = xelatex
% UTF-8. Standalone: xelatex DKA.tex from this directory.
% To include in an existing XePersian document, define \DKAInNotebook first.
% Source: Google Drive Notebook/DKA; 20 original JPEGs in filename order.
% Wording, numbers, abbreviations and handwritten corrections are retained.
% pic01.jpg = 20261012_202449096.jpg; Drive ID 1H3cCHx7gRqUgqcG33NXHTRzRZpbyYM5q
% pic02.jpg = 20261012_202457535.jpg; Drive ID 18qNq_psJHRlvvRolGJVUIkZBXGbafIBQ
% pic03.jpg = 20261012_202503118.jpg; Drive ID 1lYlsr6ujNEPvSxKnHTUY29QXrNUI7RX4
% pic04.jpg = 20261012_202509208.jpg; Drive ID 1m2_yZECKZ6T7NUDZtR1AO8s1Ldh0awQM
% pic05.jpg = 20261012_202517506.jpg; Drive ID 1N13T8UE0X3IWAdGq6tkGZ3gNkmcw573j
% pic06.jpg = 20261012_202525143.jpg; Drive ID 1DMp6smHGNwCQMKPkQgWd3kiYFOmSpKBA
% pic07.jpg = 20261012_202531477.jpg; Drive ID 11-0lVoxDvd7_pEVLuhX4XF5wT6IMoulR
% pic08.jpg = 20261012_202537579.jpg; Drive ID 1_5Q5QDtWZ22T3liQQnlcvWsGdaT5U9yB
% pic09.jpg = 20261012_202549142.jpg; Drive ID 10EUkdzTzvnYmANM4vl09YPj9pWPwWaG9
% pic10.jpg = 20261012_202556004.jpg; Drive ID 19siWUw6rgLZlbbG-InlMlY6WNHVWKQAZ
% pic11.jpg = 20261012_202602342.jpg; Drive ID 1632H_i34L6odgTIvh7JHL2OrgsuXirSe
% pic12.jpg = 20261012_202607521.jpg; Drive ID 1Dhuly05kog-nQwU1SlDkyuW5No9_Tx-a
% pic13.jpg = 20261012_202615325.jpg; Drive ID 1DsZ4owalZ_j5Xg31gEDaYvgHHK775QHW
% pic14.jpg = 20261012_202620903.jpg; Drive ID 1YfB7YmDdWhzqCW4iF7gu1iS_HPBoNXMr
% pic15.jpg = 20261012_202627830.jpg; Drive ID 1q8Gwj3yQV3Zku9H8RH-jvp60oL9JYMi4
% pic16.jpg = 20261012_202634807.jpg; Drive ID 1qict_Tyn-KpgOCGYkHrKUCfcKR7ZsVu2
% pic17.jpg = 20261012_202642433.jpg; Drive ID 11H3Bv6WH9HTQ7S23bVHUYZD1lVvgZ0YF
% pic18.jpg = 20261012_202650200.jpg; Drive ID 1Hg7yEbRO3Iokxp1C1RQ9fP_qBUjIWgDF
% pic19.jpg = 20261012_202657495.jpg; Drive ID 1VylQrJZBzBdEb8aEYsqqzSW7-4VzVEcT
% pic20.jpg = 20261012_202703708.jpg; Drive ID 1ckq8skWjV5hA5ycoFZBuein9KPoQi_lT
\ifdefined\DKAInNotebook
  \let\DKAfinish\relax
\else
  \def\DKAfinish{\end{document}}
  \documentclass[12pt,a4paper]{article}
  \usepackage[margin=22mm]{geometry}
  \usepackage{amsmath,amssymb}
  \usepackage{graphicx}
  \usepackage{array,longtable}
  \usepackage{enumitem}
  \usepackage{xepersian}
  \IfFontExistsTF{Amiri}{\settextfont{Amiri}}{\settextfont{DejaVu Sans}}
  \setlatintextfont{DejaVu Serif}
  \setlength{\parindent}{0pt}
  \setlength{\parskip}{5pt}
  \setlength{\emergencystretch}{3em}
  \begin{document}
\fi
\providecommand{\DKAunclear}[1]{\textbf{[خوانش نامطمئن: #1]}}
\providecommand{\DKApage}[2]{\clearpage\section*{تصویر #1}\lr{\texttt{#2}}\par}
\providecommand{\DKAmath}[1]{\lr{\(#1\)}}

\begin{center}
{\Large\bfseries \lr{DKA} در کودکان}\par
\lr{ISPAD}
\end{center}
\textbf{یادداشت رونویسی:} متن به ترتیب ۲۰ تصویر نوشته شده است؛ عبارت‌ها، اختصارات، اعداد و یادداشت‌های حاشیه‌ای منبع حفظ شده‌اند. فلش‌ها به صورت رابطه یا فهرست بازنویسی شده‌اند. خوانش‌های نامطمئن مشخص شده‌اند و تصاویر کامل در پیوست آمده‌اند. سربرگ چاپی تمام برگه‌ها «موضوع» و «تاریخ / /» است؛ سربرگ دست‌نویس تصاویر ۱–۳ «\lr{DKA} در کودکان» / «\lr{DKA} کودکان» و سربرگ موضوع تصاویر ۴–۲۰ خالی است. شمارهٔ دست‌نویس هر صفحه مطابق شمارهٔ تصویر است.

\DKApage{۱}{pic01.jpg}
حاشیهٔ بالای صفحه: \lr{ISPAD}.
\subsection*{کرایتریا بیوشیمیایی تشخیص \lr{DKA}}
\begin{itemize}
\item هایپرگلاسیمی: \DKAmath{BS>11\,\mathrm{mmol/L}\simeq200\,\mathrm{mg/dL}}.
\item \lr{pH} وریدی کمتر از ۷٫۳ یا بی‌کربنات کمتر از \lr{18 mmol/L}.
\item کتونمی (بتاهیدروکسی‌بوتیرات خون \DKAmath{\geq3\,\mathrm{mmol/L}}) یا کتونوری متوسط تا شدید (\DKAmath{\geq2+}).
\end{itemize}
همهٔ کودکان در زمان مراجعه علائم دیابت مثل پرادراری، پرنوشی ندارند. (در دست‌نوشته، \lr{DKA} خط خورده و «دیابت» جایگزین شده است.)

علائم \lr{DKA} غیراختصاصی هستند. سنجش سطح گلوکز با گلوکومتر باید در همهٔ کودکان با تاکی‌پنه یا استفراغ و درد شکم بدون اسهال انجام شود.

در قدم اول \lr{PALS} را انجام دهید:
\begin{itemize}
\item اندازه‌گیری سریع قند خون.
\item اندازه‌گیری کتون خون یا ادرار.
\item اندازه‌گیری الکترولیت سرم.
\item اندازه‌گیری گازهای خون.
\item ارزیابی سطح هوشیاری.
\item ۲ عدد \lr{IV} محیطی.
\end{itemize}

\DKApage{۲}{pic02.jpg}
پاسخ بالینی و آزمایشگاهی باید به طور متناوب ارزیابی شود.
\subsection*{اهداف درمان \lr{DKA}}
\begin{enumerate}
\item اصلاح دهیدریشن.
\item اصلاح اسیدوز.
\item ریورس کردن کتوز.
\item هایپراسمولالیته و سطح گلوکز تدریجاً به \lr{N/L} برگردند.
\item مانیتور از نظر عوارض حاد \lr{DKA}.
\item تشخیص زودهنگام و درمان به موقع هر رخداد آینده.
\end{enumerate}
\subsection*{مایع‌درمانی}
مایع‌درمانی باید قبل از انسولین آغاز شود.
\begin{itemize}
\item یک یا چند بولوس \lr{N/S} طی ۲۰–۳۰ دقیقه برای بازگشت جریان خون محیطی.
\item ادامهٔ مایع‌درمانی با \lr{N/S} یا \lr{H/S} به گونه‌ای که طی ۲۴–۴۸ ساعت دفیسیت مایع بیمار جبران گردد.
\end{itemize}
\subsection*{انسولین‌درمانی}
\begin{itemize}
\item اگر \DKAmath{pH>7.15}: \lr{0.05 U/kg/h} انسولین رگولار وریدی.
\item اگر \DKAmath{pH\leq7.15}: \lr{0.1 U/kg/h} رگولار وریدی.
\item حداقل ۱ ساعت پس از شروع مایع‌درمانی آغاز گردد.
\end{itemize}

\DKApage{۳}{pic03.jpg}
\subsection*{پتاسیم}
\begin{itemize}
\item اگر \DKAmath{K>5.5}: تا زمانی که \lr{Output} ادراری برقرار شود پتاسیم نمی‌دهیم.
\item مایع‌درمانی اولیه پتاسیم ندارد.
\item در ابتدا و سپس هر ۱ ساعت پتاسیم را چک می‌کنیم.
\item زمانی که \DKAmath{K<5.5}: پتاسیم وریدی می‌دهیم؛ \lr{40 mmol/L}.
\item اگر \DKAmath{K<3}: قطع انسولین و بولوس پتاسیم؛ حداکثر \lr{0.5 mEq/kg/h}؛ \lr{C.M}.
\end{itemize}
\subsection*{بی‌کربنات}
اندیکاسیون‌های بسیار محدود:
\begin{itemize}
\item هایپرکالمی تهدیدکنندهٔ حیات.
\item اسیدوز شدید (\DKAmath{pH<6.9}) + شواهد اختلال انقباض قلب.
\end{itemize}
\subsection*{علائم و نشانه‌های خطر آسیب مغزی}
\begin{itemize}
\item شروع سردرد یا شروع پس از آغاز درمان.
\item سردرد شدید یا سیر تشدیدشونده.
\item افت \lr{HR} (به علت خواب و یا اصلاح شوک نباشد).
\item \lr{MSC}.
\item شواهد فلج هر یک از اعصاب مغزی.
\item افت \DKAmath{SpO_2}.
\end{itemize}

\DKApage{۴}{pic04.jpg}
\DKAmath{BP\uparrow} اغلب در بیماران \lr{DKA} دیده می‌شود و نباید نشانهٔ خطر آسیب مغزی فرض گردد مگر شواهد دیگری نیز داشته باشیم.
\subsection*{در کودک با ریسک فاکتورهای متعدد آسیب مغزی}
\begin{itemize}
\item \DKAmath{BUN>20\,\mathrm{mg/dL}}.
\item اسیدوز شدید (\DKAmath{pH<7.1}).
\item هایپوکاپنهٔ شدید (\DKAmath{pCO_2<21}).
\item سن کمتر از ۵ سال.
\item مانیتول و سالین هایپرتونیک در \lr{bedside} بیمار باشد و دوز آن از پیش محاسبه کنید.
\item اگر وضعیت نورولوژیک به طور حاد بد شد: بلافاصله درمان.
\end{itemize}
منیج \lr{DKA} کامل نمی‌شود مگر تلاش برای شناسایی علت \lr{DKA} انجام شود.

\lr{DKA} در کودکی که کیس شناخته‌شدهٔ \lr{DM} است و \lr{illness} اخیر به عنوان علت \lr{DKA} ندارد می‌تواند به علت عدم مصرف درست انسولین باشد.

در \lr{New Onset DM}، بروز \lr{DKA} نتیجهٔ تشخیص دیرهنگام \lr{DM} است.
\subsection*{برخی داروها به طور کاذب کتون ادرار را مثبت می‌کنند}
\lr{Na-Val}، کاپتوپریل، پنی‌سیلامین، \lr{NAC}، \lr{mesna}.

\DKApage{۵}{pic05.jpg}
\subsection*{یوگلایسمیک کتواسیدوز}
\begin{itemize}
\item \lr{Starvation}.
\item رژیم با کربوهیدرات پایین و چربی بالا.
\item \lr{SGLT2 I}.
\end{itemize}
درمان یوگلایسمیک \lr{DKA} مانند \lr{DKA} معمولی است با این تفاوت که دکستروز زودتر شروع می‌شود: بلافاصله بعد از مایع‌درمانی اولیه.

سطح بی‌کربنات سرم در مواردی که به \lr{ABG} دسترسی نداریم می‌تواند در تشخیص و مانیتورینگ بیمار جایگزین \lr{pH} شود.
\subsection*{ارزیابی شدت \lr{DKA}}
\begin{latin}
\begin{itemize}
\item Mild: \(pH<7.3\) or \(S_{\mathrm{HCO}_3}<18\).
\item Moderate: \(pH<7.2\) or \(S_{\mathrm{HCO}_3}<10\).
\item Severe: \(pH<7.1\) or \(S_{\mathrm{HCO}_3}<5\).
\end{itemize}
\end{latin}
\subsection*{در شرح حال بیمار \lr{DKA}}
پرادراری، پلی‌دیپسی، تهوع، استفراغ، تاکی‌پنه، \lr{SOB}، درد شکم، ضعف، کاهش وزن، کانفیوژن، \DKAmath{LOC\downarrow}.

اگر در بدو ورود بیمار در شوک بود: هم \lr{NG}، هم اکسیژن ۱۰۰\% (نبض‌ها خیلی ضعیف بود).

\DKApage{۶}{pic06.jpg}
\subsection*{مانیتورینگ در \lr{DKA}}
\begin{itemize}
\item \lr{BS} هر ۱ ساعت.
\item تنظیم مایع هر ۱ ساعت.
\item وضعیت نورولوژیک هر ۱ ساعت (\lr{GCS}).
\item الکترولیت‌ها هر ۲ ساعت.
\item \lr{C.M} مدام.
\end{itemize}
طی درمان هیپرگلیسمی، \DKAmath{BS<300\,\mathrm{mg/dL}} (\lr{17 mmol/L}): اضافه‌کردن گلوکز ۵–۱۲٫۵\% به مایع \lr{IV}.

سیر بهبودی: \lr{Clinically well} + رفع کتواسیدوز + تحمل مایع خوراکی: شروع انسولین \lr{SC} و فاز ترنزیشن.

شک به آسیب مغزی: رد هایپوگلاسیمی؛ تأیید \lr{Cerebral Injury}.

درمان \lr{CI}:
\begin{itemize}
\item مانیتول \lr{0.5--1 g/kg} یا سالین هایپرتونیک (۳\%).
\item تنظیم مایع‌درمانی برای حفظ \lr{BP} در محدودهٔ \lr{N/L}.
\item انتقال به \lr{ICU}.
\item تصویربرداری مغز پس از پایدار شدن بیمار.
\end{itemize}

\DKApage{۷}{pic07.jpg}
در درصد اندکی از موارد \lr{DKA} در کودکان علل عفونی دارد: عمدتاً \lr{GE}.

جدا از وجود یا عدم وجود دهیدریشن، ۱۲\% بیماران در ابتدا و ۱۴\% طی روند درمان دچار \lr{HTN} هستند.

\DKAmath{BOHB\geq3} میلی‌مول در لیتر؛ اندازه‌گیری بالای سر بیمار، تأییدکنندهٔ تشخیص کتوز.

آزمایش اولیه: \lr{VBG}، \lr{BS}، الکترولیت‌ها، \lr{BUN}، \lr{Cr}.

تعیین درجهٔ دهیدریشن به طور دقیق ممکن نیست در \lr{DKA}؛ بر اساس معاینه و شرح حال.
\begin{itemize}
\item در دهیدریشن \DKAmath{\geq10\%}: نبض محیطی ضعیف، \DKAmath{BP\downarrow}، الیگوری.
\end{itemize}
در کودک بدون هوشیاری اگر رفلکس‌های محافظت‌کنندهٔ راه هوایی وجود نداشت، بیمار را انتوبه می‌کنیم.
\begin{itemize}
\item \lr{NG-Tube} با ساکشن مداوم برای جلوگیری از پنومونی آسپیراسیون.
\item تا حد امکان نباید بیمار را انتوبه کنیم.
\end{itemize}
به بیمار با شوک اکسیژن می‌دهیم.

مانیتورینگ مداوم قلبی: درجهٔ تاکی‌کاردی، آریتمی، موج \lr{T}.

\DKApage{۸}{pic08.jpg}
دسترسی شریانی معمولاً نیاز نیست مگر در کودکان \lr{Critically ill} در \lr{ICU}.
\begin{itemize}
\item تا حد ممکن از تعبیهٔ \lr{CV-line} پرهیز شود: ریسک بالای ترومبوز.
\item اگر تعبیه شد به محض اینکه شرایط بیمار اجازه داد خارج شود.
\item اگر تعبیه شد خصوصاً در کودکان بالای ۱۲ سال: \lr{LMWH}.
\item انسولین نباید از طریق \lr{CV-line} داده شود چراکه انفیوژن آن ممکن است به علت تزریق سایر داروها قطع شود.
\end{itemize}
اگر در کودک شواهد عفونت وجود داشت، کشت‌های لازم دریافت شود سپس درمان \lr{ABT} شروع گردد.

تعبیهٔ سوند معمولاً نیاز نیست مگر در بیمار غیرهوشیار یا \lr{severely ill}.
\subsection*{آزمایش‌های تکمیلی}
\begin{itemize}
\item \lr{Hb, Hct}.
\item \lr{Alb, Ca, Ph, Mg}.
\item \lr{HbA1c}.
\end{itemize}
در اکثر بیماران \lr{DKA} لکوسیتوز با شیفت به چپ وجود دارد و نباید ملاک تشخیص عفونت باشد. عفونت را بر اساس بالین و سایر ارزیابی می‌کنیم.

\DKApage{۹}{pic09.jpg}
اگر پتاسیم به صورت اورژانسی آماده نشود: \lr{EKG} می‌گیریم.
\subsection*{اندیکاسیون بستری در \lr{PICU}}
\begin{itemize}
\item \lr{Severe DKA}.
\item ریسک افزایش‌یافتهٔ آسیب مغزی.
\end{itemize}
\subsection*{اصول مایع‌درمانی و الکترولیت در \lr{DKA} کودکان}
\begin{itemize}
\item میزان تقریبی دفیسیت در کودک حدود ۷\% وزن است.
\item افت فشار، \DKAunclear{واژهٔ بعد از «افت فشار»: «شایع»؟} بیمار دچار شوک با همودینامیک ناپایدار می‌شود.
\item از آنجایی که تعیین میزان دفیسیت مایع بیمار دشوار است:
\begin{itemize}
\item در \lr{mild DKA}: ۵\% فرض می‌شود.
\item در \lr{moderate DKA}: ۷\% فرض می‌شود.
\item در \lr{severe DKA}: ۱۰\% فرض می‌شود.
\end{itemize}
\item \DKAmath{BUN\uparrow} و \DKAmath{AG\uparrow} بهترین ارتباط را با درجهٔ دهیدریشن دارند.
\item سدیم اصلاح‌شده نشانگر خوبی است که اگر هایپرگلاسیمی وجود نداشت بود. \DKAunclear{ساختار جمله در دست‌نوشته به همین صورت است}
\item مایع‌درمانی در مطالعات اخیر منجر به افزایش ریسک آسیب مغزی نشده.
\item پاتوژنز آسیب مغزی نامشخص است.
\item از سرم‌های کریستالوئیدی باید استفاده کنیم نه کلوئیدی.
\item مایع در احیا: \lr{10--20 cc/kg} طی \lr{20--30 min}؛ صرفاً \lr{N/S}.
\item دفیسیت: \lr{N/S}، \lr{H/S}، رینگر، هارتمن، \lr{Plasmalyte}.
\end{itemize}

\DKApage{۱۰}{pic10.jpg}
\begin{itemize}
\item چون به دنبال پایین آمدن سطح گلوکز خون، حجم مایع داخل عروقی جبران می‌شود: باید مایع \lr{maintenance} کافی گذاشته شود.
\item باید به سرم، پتاسیم کلراید، فسفات یا استات اضافه شود.
\item علاوه بر مایع \lr{maintenance} روزانه، حجم دفیسیت را محاسبه کنید، حجم مایع اولیه را از آن کم کنید و آن را طی ۲۴–۴۸ ساعت به بیمار بدهید.
\end{itemize}
\subsection*{معیارها و گایدهای بالینی مایع‌درمانی}
\begin{itemize}
\item وضعیت گردش خون در ارزیابی بالینی.
\item ترند سدیم سرم.
\item اگر سدیم سرم پایین باشد و با پایین آمدن گلوکز بالا نرود، باید میزان سدیم سرم تزریقی بیمار را بالا ببریم.
\item دفع ادراری نباید به حجم مایع سرم بیمار اضافه شود.
\item مگر کودک با تظاهر همزمان \lr{DKA} و \lr{HHS}.
\item مانیتورینگ دقیق \lr{I/O} بیمار ضروری است.
\item محاسبهٔ حجم مایع کودکان چاق مانند کودک عادی است.
\item از \lr{IBW} استفاده نمی‌کنیم.
\item حداکثر \lr{500 mL/h} مایع و حداکثر ۱ لیتر بولوس.
\item دادن حجم زیادی مایع با کلر بالا در حالی که کلیه دارد به جای کلر کتون دفع می‌کند: ایجاد اسیدوز متابولیک هایپرکلرمیک.
\item معمولاً بدون علامت است و خودبه‌خود رفع می‌شود.
\end{itemize}

\DKApage{۱۱}{pic11.jpg}
ممکن است \lr{BD} ماندگار و بی‌کربنات پایین به عنوان ادامهٔ کتواسیدوز تفسیر شود که اشتباه است.
\begin{itemize}
\item اندازه‌گیری سریال \lr{BOHB} یا محاسبهٔ \lr{AG} کمک‌کننده.
\item در این شرایط بهتر است نمک پتاسیم به جای \lr{KCl} استفاده شود.
\end{itemize}
\subsection*{جایگزینی پتاسیم}
\begin{itemize}
\item میزان کلی کمبود پتاسیم در کودک با \lr{DKA} در حدود \lr{3 to 6 mmol/kg}.
\item پتاسیم از طریق استفراغ و دیورز اسموتیک از بدن خارج شده است.
\item ولوم دپلیشن منجر به هایپرآلدوسترونیسم ثانویه شده است: تشدید دفع ادراری پتاسیم.
\item احتمال \DKAmath{K\downarrow} (\DKAmath{<2.5}) در کودک با سوءتغذیه بیشتر است.
\item انسولین: شیفت پتاسیم به داخل سلول.
\item اثر شبه‌آلدوسترون و افزایش دفع کلیوی پتاسیم.
\end{itemize}
اگر در ابتدا کودک هایپوکالمیک است:
\begin{itemize}
\item شروع پتاسیم همزمان با شروع مایع بولوس (قبل از انسولین).
\item در \DKAmath{K<3.5}: قطع انسولین + بولوس پتاسیم (حداکثر \lr{0.5 mmol/kg/h}) به همراه \lr{C.M}.
\item در غیر این صورت پتاسیم را همزمان با شروع انسولین شروع می‌کنیم.
\end{itemize}

\DKApage{۱۲}{pic12.jpg}
\begin{itemize}
\item اگر کودک هایپرکالمیک است، پتاسیم را تا زمان برقراری \lr{output} ادراری شروع نمی‌کنیم.
\item ارزیابی ساعتی تا \DKAmath{K<5.5}.
\item اگر پتاسیم به صورت اورژانسی چک نمی‌شود: \lr{EKG}.
\end{itemize}
\subsection*{تغییرات \lr{EKG}}
\begin{itemize}
\item \DKAmath{K\downarrow}: پرولانگ شدن \lr{PR interval}، فلت شدن موج \lr{T}، \lr{TWI}، \lr{STD}، موج \lr{U} برجسته، \lr{Long QT} (به دلیل فیوز شدن \lr{T} و \lr{U}).
\item \DKAmath{K\uparrow}: موج \lr{T} تال، پیکد، سیمیتریکال؛ \lr{Short QT interval}.
\end{itemize}
\begin{itemize}
\item دوز جایگزینی پتاسیم \lr{40 mmol/L} مایع سرم است. در ادامه دوز بر اساس اندازه‌گیری‌های بعدی تعیین می‌شود.
\item بهتر است همهٔ پتاسیم را به صورت \lr{KCl} ندهیم چون اسیدوز متابولیک هایپرکلرمیک ایجاد می‌کند.
\item پتاسیم استات / پتاسیم فسفات.
\item اگر همه را پتاسیم فسفات بدهیم: \DKAmath{Ca\downarrow}.
\item حداکثر سرعت جایگزینی پتاسیم \lr{0.5 mmol/kg/h} است.
\item اگر علی‌رغم حداکثر سرعت جایگزینی پتاسیم همچنان \DKAmath{K\downarrow} بود: سرعت انسولین را کم می‌کنیم.
\end{itemize}

\DKApage{۱۳}{pic13.jpg}
\subsection*{فسفات}
\begin{itemize}
\item طی \lr{DKA} فسفات از دست می‌رود.
\item دیورز اسموتیک.
\item شیفت فسفات داخل سلولی به بیرون به علت اسیدوز متابولیک.
\item انسولین فسفات را به داخل سلول برمی‌گرداند.
\item \DKAmath{Ph\downarrow} سرم طی درمان.
\item طی روند درمان ۵۰\% بیماران هایپوفسفاتمی می‌شوند.
\item شدت اسیدوز متابولیک و ادامهٔ درمان وریدی بدون آغاز تغذیهٔ \lr{PO} ریسک فاکتورهای اصلی هستند.
\item هایپوفسفاتمی شدید نادر است.
\end{itemize}
عوارض:
\begin{itemize}
\item افینیتی \lr{Hb} به \DKAmath{O_2} بالا می‌رود و اکسیژن‌رسانی بافتی مختل می‌گردد.
\item متابولیک انسفالوپاتی.
\item تشنج.
\item میوپاتی پروگزیمال.
\item اختلال قدرت انقباض قلب.
\item آریتمی بطنی.
\item دیسفاژی.
\item نارسایی تنفسی.
\item ایلئوس.
\item آنمی همولیتیک.
\item رابدومیولیز.
\end{itemize}
هایپوفسفاتمی (\DKAmath{<1\,\mathrm{mg/dL}} یا \DKAmath{<0.32\,\mathrm{mmol/L}}) چه علامت‌دار باشد چه نباشد باید درمان شود.
\begin{itemize}
\item انسولین به طور موقت \lr{hold} شود.
\end{itemize}

\DKApage{۱۴}{pic14.jpg}
\begin{itemize}
\item جایگزینی روتین فسفات در بیمار \lr{DKA} توصیه می‌شود.
\item طی جایگزینی فسفات سطح \lr{Mg} و \lr{Ca} پایش شود.
\end{itemize}
\subsection*{درمان با انسولین}
\begin{itemize}
\item شروع انفیوژن انسولین، یک ساعت پس از شروع مایع‌درمانی.
\item \lr{50 U} انسولین رگولار در \lr{50 mL N/S} یک محلول \lr{1 U/mL} می‌سازد.
\item باید \lr{IV} داده شود.
\item انسولین بولوس نمی‌دهیم.
\item اگر \lr{IV} نداشت، به صورت \lr{IM} نمی‌دهیم. \DKAunclear{«نمی‌دهیم» خوانده می‌شود؛ احتمال «می‌دهیم»}
\item از \lr{CV-line} استفاده نمی‌کنیم چون \lr{dead space} زیادی دارد.
\item حداقل تا رفع \lr{DKA} باید انسولین را با دوز \lr{0.05--0.1 U/kg/h} بدهیم.
\item رفع: \DKAmath{pH>7.3}، \DKAmath{HCO_3>18}، رفع \lr{AG} یا \DKAmath{BOHB<1}.
\item \lr{pH} و \lr{BOHB} هر \lr{2 h} چک شود.
\item اگر اثر انسولین مناسب باشد، \lr{BOHB} به میزان \lr{0.5 mmol/L/h} افت می‌کند.
\item در کودک زیر ۵ سال و \lr{mild DKA} انسولین را با دوز \lr{0.03} می‌دهیم.
\item سطح گلوکز باید با سرعت \lr{2--5 mmol/L per h} کاهش یابد.
\item \lr{D/W 5\%} پس از رسیدن \lr{BS} به ۲۵۰–۳۰۰ شروع شود.
\item ممکن است نیاز به ۱۰\% یا حتی ۱۲٫۵\% هم باشد.
\item در \lr{DKA mild} و \lr{moderate} بدون عارضه می‌توان از انسولین \lr{SC} سریع‌الاثر (\lr{lispro} یا \lr{Aspart}) استفاده کرد، هر ۱–۲ ساعت.
\end{itemize}

\DKApage{۱۵}{pic15.jpg}
\subsection*{انسولین \lr{SC}}
\begin{itemize}
\item شروع ۱ ساعت پس از آغاز مایع‌درمانی.
\item دوز انسولین \lr{SC}: \lr{0.15 U/kg q2h}؛ سپس \lr{0.1 U/kg q2h} اگر \lr{BS} به میزان \lr{90 mg/dL/h} افت کند.
\item تجویز انسولین رگولار به صورت \lr{SC} یک گزینهٔ دیگر است.
\item \lr{q4h}.
\item \lr{0.13--0.17 U/kg/dose}؛ افزایش یا کاهش ۱۰–۲۰\%.
\end{itemize}
\subsection*{اسیدوز}
\begin{itemize}
\item مایع و انسولین اسیدوز را ریورس می‌کنند.
\item از متابولیزه شدن کتواسیدها بی‌کربنات تولید می‌شود.
\item سرعت بالاتر مایع‌درمانی، نرمال شدن \lr{AG} را تسریع می‌کند ولی تأثیری در سرعت نرمال شدن \lr{pH} ندارد.
\item چون اسیدوز متابولیک هایپرکلرمیک رخ می‌دهد.
\item عدم رفع کامل اسیدوز مانع شروع انسولین \lr{SC} نیست چراکه علت آن اسیدوز متابولیک هایپرکلرمیک است و پروسهٔ خوش‌خیمی است.
\item \lr{Sepsis} یک علت مهم عدم رفع اسیدوز است.
\item درمان با بی‌کربنات منجر به اسیدوز پارادوکسیکال در \lr{CNS} می‌شود.
\item همچنین اصلاح سریع اسیدوز هایپوکالمی ایجاد می‌کند.
\end{itemize}
اندیکاسیون دادن بی‌کربنات:
\begin{itemize}
\item هایپرکالمی تهدیدکنندهٔ حیات.
\item اسیدوز شدید (\DKAmath{pH<6.9}) + اختلال کانترکتیلیتی قلب.
\end{itemize}

\DKApage{۱۶}{pic16.jpg}
\subsection*{شروع مایعات خوراکی و تبدیل به انسولین \lr{SC}}
\begin{itemize}
\item \lr{PO} تنها زمانی باید آغاز شود که بهبودی واقعی بالینی رخ داده باشد.
\item کتون ادرار تا ساعت‌ها پس از \lr{N/L} شدن سطح سرمی \lr{BOHB} همچنان مثبت خواهد بود.
\item در ادرار، استون و استواستات اندازه‌گیری می‌شود.
\item عدم وجود کتونوری معیار رفع \lr{DKA} نیست.
\end{itemize}
زمانی که کتواسیدوز رفع شد و تحمل \lr{PO} برقرار شد: برنامه‌ریزی به انسولین \lr{SC}. (در منبع \lr{PO} کنار انسولین خط خورده و \lr{SC} بالای آن نوشته شده است.)
\begin{itemize}
\item یک دوز انسولین بازال (طولانی‌اثر) باید شروع شود + کوتاه‌اثر (هر دو \lr{SC}).
\item بهترین زمان تبدیل به \lr{SC} دقیقاً قبل از وعدهٔ غذایی است.
\item یک رویکرد دیگر انسولین بازال \lr{SC} را زمانی می‌دهیم که بیمار همچنان در حال دریافت انسولین \lr{IV} است.
\end{itemize}
\begin{itemize}
\item برای جلوگیری از ریباند هایپرگلاسیمی، اولین تزریق انسولین کوتاه‌اثر \lr{SC} باید \lr{15--30 min} قبل از قطع انسولین \lr{IV} باشد.
\item اگر بخواهیم بازال \lr{SC} بدهیم این مدت باید طولانی‌تر باشد و انسولین \lr{IV} باید تدریجاً قطع شود.
\item نوع و دوز انسولین \lr{SC} باید بر اساس \lr{local} انتخاب شود.
\item پس از تبدیل به انسولین \lr{SC} باید چندین نوبت روزانه \lr{BS} چک شود.
\end{itemize}

\DKApage{۱۷}{pic17.jpg}
\subsection*{مانیتورینگ بیمار \lr{DKA}}
\begin{itemize}
\item هر ۱ ساعت: علائم حیاتی، \lr{GCS} و علائم خطر آسیب مغزی، \lr{I/O}، \lr{BS} مویرگی.
\item بدو بستری سپس هر ۲ ساعت: الکترولیت‌ها، قند آزمایشگاهی، \lr{BUN}، \lr{Ca}، \lr{Mg}، فسفات، \lr{VBG}، \lr{BOHB} (در صورت امکان).
\item وزن بیمار هر روز صبح اندازه‌گیری شود.
\end{itemize}
\begin{latin}
\[
\mathrm{Anion\ Gap}=\mathrm{Na}-(\mathrm{Cl}+\mathrm{HCO}_3).
\]
N/L: \(12\pm2\,\mathrm{mmol/L}\); in DKA \(\simeq20\text{--}30\).

If \(>35\), suggests DKA + lactic acidosis.
\[
\mathrm{Corrected\ Na}
=\mathrm{Measured\ Na}+
\left(1.6\times\frac{BS-100}{100}\right).
\]
\end{latin}
\DKAunclear{ضریب در فرمول سدیم اصلاح‌شده: «۱٫۶»؛ احتمال خوانش «۱٫۴»}
\begin{latin}
\[
\mathrm{Effective\ Osmolality}=(2\times\mathrm{Na})+\mathrm{Glu}
\qquad \mathrm{mOsm/kg}.
\]
N/L is \(275\text{--}275\).
\end{latin}
\textbf{یادداشت رونویسی:} بازهٔ اسمولالیته در تصویر دقیقاً \lr{275--275} نوشته شده است.
\subsection*{عوارض \lr{DKA} در کودکان}
\begin{itemize}
\item شایع‌ترین علت مرگ در \lr{DKA} آسیب مغزی است.
\item در توسعه‌یافته کمتر از ۱\%؛ در حال توسعه ۳–۱۳\% نرخ مرگ.
\item می‌تواند همراه با اختلال حافظه باشد.
\item می‌تواند منجر به اختلال در هوش کلی گردد.
\end{itemize}

\DKApage{۱۸}{pic18.jpg}
\subsection*{\lr{AKI} و \lr{RTD} (آسیب رنال توبولار)}
\begin{itemize}
\item در ۴۳–۶۴\% کودکان بستری.
\item شدت \lr{AKI} وابسته به میزان کمبود مایع و شدت اسیدوز است.
\end{itemize}
\subsection*{آسیب مغزی (\lr{CI})}
\begin{itemize}
\item شیوع ۰٫۵–۰٫۹\% بیماران \lr{DKA}.
\item مرگ‌ومیر ۲۴\%.
\item در حدود ۱۵\% بیماران \lr{DKA}، \DKAmath{GCS<14} دیده می‌شود که معمولاً در زمینهٔ ادم مغزی خفیف است.
\item در گذشته فکر می‌کردند علت اصلی آن مایع‌درمانی سریع است ولی در حال حاضر هایپوپرفیوژن مغز و وضعیت هایپراینفلامیتوری ناشی از \lr{DKA} علل اصلی در نظر گرفته می‌شوند.
\end{itemize}
درجهٔ ادم مغزی بیمار \lr{DKA} با:
\begin{itemize}
\item این موارد مرتبط است: شدت دهیدریشن و هایپرونتیلیشن در بدو ورود.
\item این موارد مرتبط نیست: اسمولالیتهٔ اولیه یا تغییرات اسموتیک طی درمان.
\end{itemize}
\begin{itemize}
\item در موارد شدید \lr{DKA} اختلال در \lr{BBB} رخ می‌دهد.
\item \lr{CI} در کودکان کم‌سن‌تر شایع‌تر است.
\item و در موارد \lr{New onset DM}.
\item معمولاً ظرف \lr{12 h} اول شروع درمان رخ می‌دهد؛ حداکثر \lr{48 h}.
\item برای تشخیص نیاز به تصویربرداری مغز نیست.
\item پس از تشخیص باید بلافاصله درمان آغاز شود.
\end{itemize}

\DKApage{۱۹}{pic19.jpg}
\subsection*{کرایتریای تشخیص \lr{CI}}
\begin{itemize}
\item ۱ معیار تشخیصی.
\item ۲ معیار ماژور.
\item ۱ معیار ماژور + ۲ معیار مینور.
\end{itemize}
\lr{Sens 92\%}؛ \lr{FP 4\%}.
\subsection*{معیارهای تشخیصی}
\begin{latin}
\begin{itemize}
\item Abnormal motor or verbal response to pain.
\item Decorticate or decerebrate posture.
\item Cranial Nerve Palsy (specially III, IV, VI).
\item Abnormal Neurologic Respiratory pattern (grunting, Tachypnea, Cheyne--Stokes respiration, Apneusis).
\end{itemize}
\end{latin}
\subsection*{معیارهای ماژور}
\begin{latin}
\begin{itemize}
\item MSC.
\item Sustained HR deceleration (\(\downarrow\) more than 20 beats/min) not attributed to shock management or sleep.
\item Age-inappropriate incontinence.
\end{itemize}
\end{latin}
\subsection*{معیارهای مینور}
\begin{latin}
\begin{itemize}
\item Vomiting.
\item Headache.
\item Lethargy.
\item \(DBP>90\,\mathrm{mmHg}\).
\item Age \(<5\).
\end{itemize}
\end{latin}

\DKApage{۲۰}{pic20.jpg}
\subsection*{درمان \lr{CI}}
\begin{itemize}
\item مانیتول: \lr{0.5--1 g/kg IV over 10--15 min}.
\item شروع اثر ظرف \lr{15 min}؛ طول مدت اثر \lr{120 min}.
\item در صورت نیاز تکرار \lr{30 min} بعد از دوز اول.
\item سالین هایپرتونیک (۳\%): \lr{2.5--5 mL/kg over 10--15 min}.
\item اگر مانیتول نبود یا اثر مانیتول کافی نبود، بعد \lr{15--30 min}.
\item در صورت احتمال نارسایی تنفسی ممکن است نیاز به انتوبیشن باشد.
\item برای بیمار انتوبه، سطح \DKAmath{PCO_2} باید در اندازهٔ قابل انتظار برای سطح اسیدوز متابولیک باشد.
\item از هایپوکاپنهٔ بیشتر از آن باید پرهیز شود: تشدید \lr{CI}.
\item پس از آغاز درمان می‌توانیم تصویربرداری انجام دهیم.
\item بررسی از نظر \lr{ICH}: نوروسرجری.
\item بررسی از نظر ترومبوز خصوصاً اگر \lr{END} داشته باشد: ضدانعقاد.
\end{itemize}

\clearpage
\section*{پیوست: تصاویر کامل منبع}
تصاویر زیر متن کامل دست‌نویس، حاشیه‌ها، سربرگ‌ها، شماره‌ها، فلش‌ها و اصلاحات اصلی را حفظ می‌کنند.
% Resolve image paths when compiled from this directory or the repository root.
\IfFileExists{images/pic01.jpg}{\def\DKAimagebase{images/}}{%
  \def\DKAimagebase{pediatrics/ped-endocrinology/DKA/images/}}
\newcommand{\DKAoriginal}[2]{%
  \clearpage\subsection*{اصل تصویر #1: \lr{\texttt{pic#2.jpg}}}%
  \begin{center}%
  \includegraphics[width=\linewidth,height=0.87\textheight,keepaspectratio]{\DKAimagebase pic#2.jpg}%
  \end{center}}
\DKAoriginal{۱}{01}
\DKAoriginal{۲}{02}
\DKAoriginal{۳}{03}
\DKAoriginal{۴}{04}
\DKAoriginal{۵}{05}
\DKAoriginal{۶}{06}
\DKAoriginal{۷}{07}
\DKAoriginal{۸}{08}
\DKAoriginal{۹}{09}
\DKAoriginal{۱۰}{10}
\DKAoriginal{۱۱}{11}
\DKAoriginal{۱۲}{12}
\DKAoriginal{۱۳}{13}
\DKAoriginal{۱۴}{14}
\DKAoriginal{۱۵}{15}
\DKAoriginal{۱۶}{16}
\DKAoriginal{۱۷}{17}
\DKAoriginal{۱۸}{18}
\DKAoriginal{۱۹}{19}
\DKAoriginal{۲۰}{20}
\DKAfinish
